\documentclass[11pt,a4paper]{article}

% --- Packages for Academic Layout & Formatting ---
\usepackage[utf8]{inputenc}
\usepackage[margin=1in]{geometry}
\usepackage{amsmath}
\usepackage{amsfonts}
\usepackage{amssymb}
\usepackage{booktabs}
\usepackage{xcolor}
\usepackage{microtype}
\usepackage{indentfirst}
\usepackage{listings}
\usepackage{graphicx}
\usepackage[most]{tcolorbox}
\usepackage{booktabs}   
\usepackage{tabularx}   
\usepackage{longtable}
\usepackage{xltabular}
\keepXColumns

% --- Professional Color Palette Selection ---
\definecolor{DeepNavy}{HTML}{1B365D}
\definecolor{SteelBlue}{HTML}{5C768D}
\definecolor{Charcoal}{HTML}{222222}
\definecolor{LightTint}{HTML}{F0F4F8}

% --- Layout Configurations ---
\usepackage{titlesec}
\titleformat{\section}{\color{DeepNavy}\normalfont\Large\bfseries}{\thesection}{1em}{}[{\color{SteelBlue}\titlerule[0.8pt]}]
\titleformat{\subsection}{\color{SteelBlue}\normalfont\large\bfseries}{\thesubsection}{1em}{}

\usepackage{setspace}
\setstretch{1.15}

\begin{document}

% =========================================================================
% TITLE PAGE
% =========================================================================
\begin{center}
    {\LARGE \textbf{\color{DeepNavy}A Unified Defense-in-Depth Framework for Mitigating LLM Hallucinations: Integrating Internal Probes with External Symbolic Verification}}\\[1.5cm]
    {\Large Master of Technology (M.Tech) Thesis Architecture Document}\\[1cm]
\end{center}

\pagenumbering{arabic}

% =========================================================================
% 1. ABSTRACT
% =========================================================================
\section{Abstract}
The deployment of Large Language Models (LLMs) into safety-critical domains is restricted by their propensity to generate hallucinations. Current interventions are highly fragmented: pretraining methods fail to adjust to live context shifts, naive Retrieval-Augmented Generation (RAG) is prone to noise injection, and multi-agent debate protocols incur prohibitive computational latency. This research proposes a pure Python, asynchronous defense-in-depth framework designed to reconcile detection precision with runtime economic efficiency. 

The framework functions via a multi-tiered architecture: (1) an internal triage layer utilizing Semantic Entropy Probes (SEPs) to instantaneously measure model uncertainty; (2) a targeted external verification track running Hypothesis-Testing RAG coupled with Natural Language Inference (NLI) checks; and (3) a structured logical reasoning track utilizing the HalluClean paradigm to decouple plan formulation from raw generation. Empirical evaluations will demonstrate improvements in operational cost-effectiveness and reductions in final hallucinatory outputs.

% =========================================================================
% 2. PROPOSED METHODOLOGY & SYSTEM ARCHITECTURE
% =========================================================================
\section{Proposed Methodology \& Pure Python Architecture}
The proposed framework establishes an integrated, defense-in-depth pipeline that transitions from fast, lightweight internal validation to detailed external verification based on real-time query risk. The entire system is constructed using an asynchronous, pure Python microservice architecture (FastAPI) to ensure zero network latency during tensor activation extraction.

\subsection{Layer 1: Internal Triage via Semantic Entropy Probes (SEPs)}
When a user prompt enters the system, the base autoregressive model (hosted locally via Ollama or Hugging Face) initiates a draft generation. The framework probes the final hidden-state activations of the tokens during the forward pass. Activations are extracted from the later one-third of the target transformer architecture to maximize semantic abstraction.

Let the hidden state vector extracted from layer $l$ at token position $t$ be represented as $h_t^{(l)} \in \mathbb{R}^d$. The linear probe parameterizes a lightweight logistical estimator:
\begin{equation}
P(\text{Hallucination} \mid h_t^{(l)}) = \sigma\left(W_{\text{SEP}} \cdot h_t^{(l)} + b_{\text{SEP}}\right)
\end{equation}
Where $\sigma$ represents the standard sigmoid activation function, and $W_{\text{SEP}} \in \mathbb{R}^{1 \times d}$ is the trained weight matrix. The semantic entropy is computed, and a predefined safety threshold ($T_{safe}$) is applied. If $H_{sem} < T_{safe}$, the token stream is marked as secure and delivered immediately, bypassing expensive validation.

\subsection{Layer 2: Hypothesis-Testing RAG with NLI Verification}
If $H_{sem} \ge T_{safe}$ and the uncertainty is knowledge-based, the system initiates a structured hypothesis-testing framework using asynchronous Python workers.
Rather than appending raw retrieved paragraphs directly to the prompt, the system extracts the model's draft claim into independent Atomic Propositions (APs). 
Retrieved source documents are processed by a fine-tuned Natural Language Inference (NLI) model (e.g., DeBERTa-v3), evaluating the relationship according to three strict states:
\begin{itemize}
    \item \textbf{Entailment:} The claim is approved for generation compilation.
    \item \textbf{Contradiction:} The AP is rejected, and negative constraints are injected into the refinement loop to force omission of the false entity.
    \item \textbf{Neutral:} The external evidence is inconclusive, triggering an expanded vector search.
\end{itemize}

\subsection{Layer 3: Structured Reasoning via the HalluClean Paradigm}
When the high uncertainty score is identified as a logical reasoning flaw, the framework routes the query into an asynchronous directed acyclic graph (DAG) pipeline.
This protocol splits the reasoning sequence into four distinct stages managed by Python state executors:
\begin{itemize}
    \item \textbf{Stage 1 (Task-Oriented Planning):} Generates a structured, symbolic execution graph.
    \item \textbf{Stage 2 (Plan-Guided Reasoning):} Executes each step sequentially within local context boundaries.
    \item \textbf{Stage 3 (Final Judgment):} Compares the synthesized conclusion against original constraints.
    \item \textbf{Stage 4 (Content Refinement):} Rewrites the verified logical path into natural language.
\end{itemize}

% =========================================================================
% 3. EVALUATION METRICS AND BENCHMARKING
% =========================================================================
\section{Evaluation Metrics and Benchmarking}
To validate the architectural efficiency, the system includes a dedicated Python evaluation harness utilizing datasets such as HaluEval 2.0 and TruthfulQA. The evaluation tracks:
\begin{itemize}
    \item \textbf{Factual Precision:} The reduction in absolute hallucination rates compared to a zero-shot baseline.
    \item \textbf{Latency Delta ($T_{total}$):} The wall-clock execution time reduction achieved by bypassing external verification for low-risk queries via Layer 1 triage.
    \item \textbf{Token Expenditure:} The aggregate reduction in autoregressive token generation across the pipeline.
\end{itemize}

\end{document}